\section{Scope, evidence, and the main recommendation}\label{sec:scope}

The target is not a generic theorem prover. It is a Lean-native synthesizer
that accepts a type and an optional behavioral contract, produces inhabitants,
and checks the resulting program and evidence against a trust profile.
That purpose makes the interaction between data construction, proof search,
and evaluation more important than any one technique in isolation.

The repository was inspected at commit
\begin{center}\small\texttt{784664ff34e819422510a4d470ab07fe48bb736b}.\end{center}
The main evidence is the complete set of research sections in
\code{docs/proposals/11-further-improvements}, the implementation notes, and
selected source modules for construction search, acceptance, and transactional
state. The proposal's benchmark and profile figures
refer to an earlier measured commit, \code{33cec8b}, using Lean 4.34.0 on one
Windows workstation. Those figures are treated as reported results, not as
measurements repeated for this article \cite{repo,proposal,implementation}.

The proposal reports complete outcome-category agreement on its scored Leant
suites, while identifying thirteen unscored Church-operation stretch cases
and additional failures just outside the corpus. That is evidence that the
existing search is useful and that the saturated baseline is insufficient
for choosing future algorithms. It is not evidence that all those failures
have the same cause or that a rewrite will improve them. Likewise, the reported
45\% residual-evaluation cost and 1\% snapshot cost come from a particular
contract-heavy profile, not a universal cost model \cite{proposal}.

The central recommendation is to retain the current kernel-based acceptance
boundary while making three distinctions architectural rather than informal:
\emph{what may be proposed}, \emph{what may be pruned}, and \emph{what may be
claimed}. An arbitrary carrier guess can be useful. An arbitrary carrier
refutation can silently remove the only solution. A checked program is useful
without being the smallest program, the intended program, or a proof of
nonexistence of alternatives. A stronger engine needs to represent these
differences explicitly.

\subsection{Kinds of conclusion used in this article}
A \emph{published result} is attributed to a primary paper or an official
implementation. A \emph{derived claim} is accompanied by an argument and
its assumptions. An \emph{engineering recommendation} is a proposed change,
not a measured improvement. A \emph{bounded guarantee} names the language,
resources, and decision procedures within which it holds. An unanswered
empirical question is accompanied by an experiment rather than a guessed
number.

The mathematical propositions below are ordinary written proofs, not Lean
formalizations. The accompanying Python program checks finite examples and
small exhaustive models; it neither executes Leant2 nor validates Lean's
kernel. The distinction matters especially for higher-order unification,
partial evaluation, and performance: the proposed adaptations have not been
implemented or benchmarked here.

\section{Corrections that change the plan}\label{sec:corrections}

\paragraph{Higher-order examples are not an unfilled Smyth feature.}
The Smyth paper explicitly distinguishes a first-order presentation of its
core calculus from an implementation supporting higher-order function examples.
It also removes the requirement that users supply trace-complete examples.
Thus the proposal's statements that Smyth's unevaluation is simply first-order
and that it shares the trace-completeness requirement should be revised.
Full dependent-type support remains a separate adaptation \cite{smyth}.

\paragraph{Auxiliary-state invention has substantial precedents.}
AutoLifter synthesizes auxiliary computations and combiners for
algorithmic lifting; Eguchi, Kobayashi, and Tsukada synthesize auxiliary
function types in templates; Para uses paramorphism-based synthesis
templates. None is a drop-in solution to the exact Lean problem, but together
they invalidate a broad conclusion that the literature does not address
auxiliary state or types \cite{autolifter,eguchi,para}.

\paragraph{Type classes already occur in type-guided abstraction refinement.}
Hoogle+'s Section~2.5 explicitly uses dictionary passing. The unresolved part
is the interaction with Lean's dependent signatures, universes, local
instances, and computational dictionaries, not the basic possibility of
representing class constraints in a synthesis net \cite{hoogleplus}.

\paragraph{Indexed induction is partly an integration problem.}
Canonical explicitly supports indexed inductive types and generalized
induction. In Lean 4.34.0, \code{Lean.Meta.mkRecursorAppPrefix} constructs a
motive by abstraction, and \code{MVarId.induction} already reverts indices and
the major premise with their dependencies. Leant2 must still select useful
extra generalizations, but it should not reimplement this existing machinery
as though it were absent \cite{canonical,leaninduction}.

\paragraph{Several proposed guarantees need stronger hypotheses.}
One unit per unfinished goal can overestimate remaining search cost.
Different partial observation rows need not describe different states.
Values already enumerated for a carrier are not an upper bound on its possible
values. Ranking after parallel lanes stop does not restore candidates that a
slower lane never generated. These are mathematical issues, not tuning
questions; counterexamples and corrected guarantees are given below.

\paragraph{Two concrete examples should be replaced.}
Ordinary list reversal is a right fold into \emph{the result type}:
\[
 \mathsf{reverse}(xs)=\foldr\,(\lambda a\,r.\ r\mathbin{+\!+}[a])\,[]\,xs.
\]
The issue is the availability and cost of the step, not the nonexistence of
such a carrier. Also, for Lean's usual natural-number addition, $n+0$ reduces
definitionally to $n$; use $0+n$ with symbolic $n$, or a commuted addition,
to illustrate a genuine propositional index transport. These observations
are used in the carrier and transport designs, respectively \cite{leanprelude}.

\section{The semantic contract of the search}\label{sec:semantic}

Fix an imported environment $E$, a local telescope $\Gamma$, a target type
$T$, and a contract $C:T\to\mathsf{Prop}$. Open terms are interpreted with a
metavariable state $\sigma$, including universe assignments and delayed
constraints. Write $\Comp_B(p,\sigma)$ for the well-typed completions admitted
by an explicitly specified bounded grammar $B$. Omitting $B$ means all
admissible completions under the chosen language and trust policy.

\begin{definition}[Three independent obligations]
\emph{Acceptance soundness} requires an accepted $t$ and $\pi$ to check as
$t:T$ and $\pi:C(t)$ in the intended environment. \emph{Conservative pruning}
requires
\[
 \mathsf{reject}(p,\sigma)\Longrightarrow
 \forall t\in\Comp(p,\sigma),\ \neg C(t).
 \tag{1}\label{eq:pruning}
\]
\emph{Optimality} requires that the returned answer minimize a specified
objective over a specified candidate set. Neither of the first two
obligations implies the third.
\end{definition}

A bounded rejection may replace $\Comp$ by $\Comp_B$, but its receipt must then
name $B$. A failed tactic invocation does not establish \eqref{eq:pruning}.
A timeout does not establish bounded exhaustion. A passing test establishes
that test, not an untested universal contract. These distinctions should
survive all the way to the result algebra and user interface.

The inspected gate synchronously checks declarations, generalizes universes
jointly over the program and proof, and audits their transitive axiom closure.
Its proof type is supplied by its caller. This is a reasonable internal API
under a trusted caller invariant, but an externally exposed synthesis service
should take the original query and reconstruct $C(t)$ inside the gate. Checking
an arbitrary supplied proof type would not by itself establish the requested
contract. This is an API-hardening recommendation, not a demonstrated
end-to-end defect in the current callers \cite{gatecode}.

Three result forms should therefore be distinct: a certified solution;
a checked impossibility proof or a genuinely exhausted, named finite search;
and an incomplete search with unrefuted survivors. Evidence supporting a
finite contract is still a fully valid proof of \emph{that finite contract}.
It simply does not identify the unstated intended algorithm.
